\chapter{Resumen}


\textbf{Resumen:}
El presente trabajo propone el desarrollo de un sistema de localización acústica y mejora del habla mediante un arreglo multimicrófono inspirado en los mecanismos de detección vibracional de los escorpiones. La motivación surge de la dificultad para aislar señales de voz en entornos acústicamente adversos, donde la presencia de múltiples fuentes sonoras y ruido degrada su inteligibilidad. El sistema integrará módulos de adquisición, localización acústica y procesamiento espacial mediante beamforming, con el objetivo de resaltar la señal de interés y atenuar el ruido de fondo. Como parte del desarrollo, se plantea adaptar principios biomiméticos asociados a la distribución espacial y heterogeneidad geométrica del sistema mecanosensorial del escorpión hacia un arreglo acústico no uniforme. El desempeño del sistema será evaluado mediante métricas objetivas como STOI, PESQ y mejora en SNR, además de compararse contra un arreglo lineal uniforme convencional procesado mediante delay-and-sum beamforming, con el fin de determinar si el enfoque propuesto ofrece ventajas prácticas en tareas de localización acústica y mejora de inteligibilidad del habla en entornos ruidosos.



\textbf{Palabras Clave:} Sistema de localización acústica, arreglo multimicrófono, Slit sensilla, escorpiones, mejora del habla, beamforming \\



%%%%%%fin del archivo

\endinput 